%% kapa: exam=2 round=37 year=2026 period=1 subject=appraisal_practice subject_ko=감정평가실무 seq=1 number=1 points=40
%% source: Q-Net 감정평가사 기출문제 (cst00309) artlSeq=5268128 file=raw/1008/5268128/26년 제37회 감정평가사 2차 (1교시_감정평가실무).pdf
%% conversion: kordoc 4.24.51 -> tools/gen2.py (auto); review=proofread-sonnet; flags=none
\begin{longquestion}{【문제 1】}
감정평가사 甲은 서울 A 지방법원으로부터 가로주택정비사업구역 내에 소재하는 분양신청을 하지 아니한 토지에 대하여 감정평가를 의뢰받았다. 제시된 자료와 관련법규 및 감정평가이론 등을 활용하여 다음 물음에 답하시오. (40점) \par
\medskip\noindent 물음 1) 공시지가기준법에 의한 시산가액을 구하시오. (10점) \par
\medskip\noindent 물음 2) 개발법(수입현가에서 비용현가 차감)에 의한 시산가액을 구하시오. (10점) \par
\medskip\noindent 물음 3) 제시된 <자료 1> 감정사항을 관련지어 본건에 적용한 공시지가기준법 및 개발법의 장ㆍ단점을 기술하시오. (10점) \par
\medskip\noindent 물음 4) 시산가액을 검토ㆍ조정하여 최종 감정평가액을 결정하시오. (10점) \par
< 자료 1 > 감정신청서 주요 사항 \par
1. 사건번호 등 \par
- 사건: 2025가단1234 소유권이전등기 \par
- 원고: M동 일대 가로주택정비사업조합 \par
- 피고: 김○○ \par
2. 감정의 목적: 원고와 피고 사이에 2026. 2. 20. 성립한 것으로 의제된 매매계약에 따라 원고가 피고에게 지급해야 하는 매매대금을 객관적으로 확정 \par
3. 감정목적물: 서울특별시 M구 M동 567-12 도로 108m\textsuperscript{2} \par
4. 감정사항: 위 감정목적물에 대하여 기준시점(2026. 2. 20.) 현재 M동 일대 가로주택정비사업의 시행으로 발생할 것으로 예상되는 개발이익을 포함한 시가 \par
< 자료 2 > 대상토지의 개황 \par
1. 개요 \par
{\small\centering\footnotesize\setlength{\tabcolsep}{3pt}
\begin{tabular}{|c|c|c|c|c|c|c|c|c|}
\hline
\makecell{기호} & \makecell{소재지} & \makecell{면적\newline{}(m\textsuperscript{2})} & \makecell{지목} & \makecell{이용상황} & \makecell{용도\newline{}지역} & \makecell{도로\newline{}교통} & \makecell{형상\newline{}지세} & \makecell{개별공시지가\newline{}(원/m\textsuperscript{2})} \\
\hline
\makecell{1} & \makecell{M동\newline{}567-12} & \makecell{108.0} & \makecell{도로} & \makecell{도로\newline{}등} & \makecell{2종일주} & \makecell{세로(가)} & \makecell{부정형\newline{}완경사} & \makecell{1,810,000} \\
\hline
\end{tabular}\par}

※ 개별공시지가 공시기준일은 2025. 1. 1. \par
2. 지리적 위치 및 주위환경 \par
본건은 서울특별시 M구 H동 소재 도시철도 ‘M역’ 북측 원거리에 소재하고, 인근은 본건 토지를 포함하여 가로주택정비사업이 추진중인 간선도로 후면에 위치한 기존 주택지대임 \par
3. 교통상황 \par
본건 북동측 인근에 마을버스정류장, 동측 인근에 시내버스정류장 그리고 남측 원거리에 도시철도 ‘M역’이 각각 소재하고 있어 대중교통사정은 보통임 \par
4. 형상 및 이용상황 등 \par
부정형 완경사로, 인접한 5개 단독주택의 출입을 위한 ‘도로’로 이용되고 있음 \par
5. 접면도로 상태 \par
지적도상 본건 토지 남동측으로 노폭 약 6m 포장도로에 접하고 있음 \par
6. 토지이용계획(주요)사항 \par
제2종일반주거지역, 소규모주택정비 관리지역(2024. 5. 10.)<빈집 및 소규모주택 정비에 관한 특례법>, 토지거래계약에관한허가구역(2025. 8. 14.) \par
< 자료 3 > 표준지 공시지가 등 \par
1. 표준지공시지가 내역 \par
[공시기준일: 2026. 1. 1.] \par
{\small\centering\footnotesize\setlength{\tabcolsep}{3pt}
\begin{tabular}{|c|c|c|c|c|c|c|c|c|}
\hline
\makecell{기호} & \makecell{소재지} & \makecell{면적\newline{}(m\textsuperscript{2})} & \makecell{지목} & \makecell{이용\newline{}상황} & \makecell{용도\newline{}지역} & \makecell{도로\newline{}접면} & \makecell{형상\newline{}지세} & \makecell{공시지가\newline{}(원/m\textsuperscript{2})} \\
\hline
\makecell{A} & \makecell{M동\newline{}458-2} & \makecell{129.1} & \makecell{대} & \makecell{단독\newline{}주택} & \makecell{2종\newline{}일주} & \makecell{세로(가)} & \makecell{사다리\newline{}평지} & \makecell{5,270,000} \\
\hline
\makecell{B} & \makecell{M동\newline{}501-9} & \makecell{118.0} & \makecell{대} & \makecell{단독\newline{}주택} & \makecell{2종\newline{}일주} & \makecell{세로(가)} & \makecell{정방형\newline{}평지} & \makecell{5,560,000} \\
\hline
\makecell{C} & \makecell{M동\newline{}568-2} & \makecell{197.8} & \makecell{대} & \makecell{다세대} & \makecell{2종\newline{}일주} & \makecell{소로\newline{}한면} & \makecell{세장형\newline{}평지} & \makecell{6,690,000} \\
\hline
\makecell{D} & \makecell{M동\newline{}570-1} & \makecell{135.4} & \makecell{대} & \makecell{단독\newline{}주택} & \makecell{2종\newline{}일주} & \makecell{세로(가)} & \makecell{정방형\newline{}평지} & \makecell{5,610,000} \\
\hline
\makecell{E} & \makecell{H동\newline{}450-1} & \makecell{187.2} & \makecell{대} & \makecell{다세대} & \makecell{2종\newline{}일주} & \makecell{소로\newline{}각지} & \makecell{사다리\newline{}평지} & \makecell{7,300,000} \\
\hline
\makecell{F} & \makecell{H동\newline{}451-4} & \makecell{165.0} & \makecell{대} & \makecell{단독\newline{}주택} & \makecell{2종\newline{}일주} & \makecell{세각(가)} & \makecell{정방형\newline{}평지} & \makecell{6,490,000} \\
\hline
\end{tabular}\par}

2. 공시지가 표준지 토지이용계획(주요)사항 \par
{\small\centering\setlength{\tabcolsep}{3pt}
\begin{tabular}{|c|>{\centering\arraybackslash}p{0.47\linewidth}|>{\centering\arraybackslash}p{0.40\linewidth}|}
\hline
기호 & 토지이용계획(주요)사항 & 비고 \\
\hline
A & - & - \\
\hline
B & - & - \\
\hline
C & 소규모주택정비 관리지역(2024. 5.\newline{}10.)<빈집 및 소규모주택정비에 관한 특례법>, 토지거래계약에관한허가구역(2025. 8. 14.) & 본건 ‘M동 일대 가로주택정비사업’\newline{}구역 내 \\
\hline
D & 위와 같음 & 본건 ‘M동 일대 가로주택정비사업’\newline{}구역 내 \\
\hline
E & 위와 같음 & ‘H동 일대 가로주택정비사업구역’\newline{}내(정비계획 내용 및 추진 일정 등에서 본건 정비사업보다 우세) \\
\hline
F & 위와 같음 & ‘H동 일대 가로주택정비사업구역’\newline{}내(정비계획 내용 및 추진 일정 등에서 본건 정비사업보다 우세) \\
\hline
\end{tabular}\par}

3. 표준지 공시지가수준 검토 \par
- ‘H동 일대 가로주택정비사업’구역은 지역의 중심축이자 제반 생활편익시설이 모여 있는 도시철도 ‘M역’ 동측 인근에 소재하고, 본건이 속한 ‘M동 일대 가로주택정비사업’구역은 도시철도 ‘M역’ 북측 원거리에 소재하고 있음 \par
- 이에 본건 사업구역이 전반적으로 ‘H동 일대 가로주택정비사업’구역에 비하여 지역요인이 열세함 \par
- 이러한 내용은 표준지 공시지가수준에도 반영되어 있는 것으로 조사된바, 지역요인 비교 시에 양 가로주택정비사업구역 내에 소재하는 표준지를 모두 활용하기로 함 \par
< 자료 4 > 지가변동률(서울특별시 M구 주거지역) \par
{\small\centering
\begin{tabular}{|c|c|c|}
\hline
\makecell{기 간} & \makecell{지가변동률(\%)} & \makecell{비 고} \\
\hline
\makecell{2026. 01. 01. ∼ 2026. 02. 20.} & \makecell{0.336} & \makecell{-} \\
\hline
\makecell{2025. 12. 23. ∼ 2026. 02. 20.} & \makecell{0.341} & \makecell{-} \\
\hline
\makecell{2026. 01. 07. ∼ 2026. 02. 20.} & \makecell{0.330} & \makecell{-} \\
\hline
\makecell{2025. 03. 15. ∼ 2026. 02. 20.} & \makecell{1.340} & \makecell{-} \\
\hline
\makecell{2025. 11. 02. ∼ 2026. 02. 20.} & \makecell{0.352} & \makecell{-} \\
\hline
\makecell{2026. 02. 20. ∼ 2026. 03. 23.} & \makecell{0.695} & \makecell{2026년 3월분 미고시로\newline{}2월분 연장 적용} \\
\hline
\end{tabular}\par}

< 자료 5 > 가치형성요인 비교 자료 \par
1. 지역요인 비교 \par
서울특별시 M구 ‘M동 일대 가로주택정비사업’ 구역과 ‘H동 일대 가로주택정비사업’ 구역의 지역요인 비교는 <자료 3> 표준지 공시지가를 활용하여 비교치를 산출함 \par
2. 개별요인 관련 자료 \par
-M동(주거지역) \par
{\small\centering
\begin{tabular}{|c|c|c|c|c|}
\hline
\makecell{이용상황} & \makecell{주거용} & \makecell{상업업무} & \makecell{주상용} & \makecell{공업용} \\
\hline
\makecell{주거용} & \makecell{1.00} & \makecell{1.34} & \makecell{1.18} & \makecell{1.08} \\
\hline
\makecell{상업업무} & \makecell{0.75} & \makecell{1.00} & \makecell{0.88} & \makecell{0.81} \\
\hline
\makecell{주상용} & \makecell{0.85} & \makecell{1.14} & \makecell{1.00} & \makecell{0.92} \\
\hline
\makecell{공업용} & \makecell{0.93} & \makecell{1.24} & \makecell{1.09} & \makecell{1.00} \\
\hline
\makecell{도로접면} & \makecell{소로한면} & \makecell{소로각지} & \makecell{세로(가)} & \makecell{세각(가)} \\
\hline
\makecell{소로한면} & \makecell{1.00} & \makecell{1.04} & \makecell{0.90} & \makecell{0.99} \\
\hline
\makecell{소로각지} & \makecell{0.97} & \makecell{1.00} & \makecell{0.87} & \makecell{0.95} \\
\hline
\makecell{세로(가)} & \makecell{1.11} & \makecell{1.15} & \makecell{1.00} & \makecell{1.09} \\
\hline
\makecell{세각(가)} & \makecell{1.01} & \makecell{1.05} & \makecell{0.91} & \makecell{1.00} \\
\hline
\end{tabular}\par}

{\small\centering
\begin{tabular}{|c|c|c|c|c|c|}
\hline
\makecell{고저} & \makecell{저지} & \makecell{평지} & \makecell{완경사} & \makecell{급경사} & \makecell{고지} \\
\hline
\makecell{저지} & \makecell{1.00} & \makecell{1.07} & \makecell{0.99} & \makecell{0.91} & \makecell{0.89} \\
\hline
\makecell{평지} & \makecell{0.93} & \makecell{1.00} & \makecell{0.93} & \makecell{0.85} & \makecell{0.83} \\
\hline
\makecell{완경사} & \makecell{1.01} & \makecell{1.08} & \makecell{1.00} & \makecell{0.92} & \makecell{0.90} \\
\hline
\makecell{급경사} & \makecell{1.10} & \makecell{1.18} & \makecell{1.09} & \makecell{1.00} & \makecell{0.98} \\
\hline
\makecell{고지} & \makecell{1.12} & \makecell{1.20} & \makecell{1.11} & \makecell{1.02} & \makecell{1.00} \\
\hline
\end{tabular}\par}

{\small\centering
\begin{tabular}{|c|c|c|c|c|c|}
\hline
\makecell{형상(주거용)} & \makecell{정방형} & \makecell{장방형} & \makecell{사다리} & \makecell{부정형} & \makecell{자루형} \\
\hline
\makecell{정방형} & \makecell{1.00} & \makecell{1.00} & \makecell{0.99} & \makecell{0.94} & \makecell{0.88} \\
\hline
\makecell{장방형} & \makecell{1.00} & \makecell{1.00} & \makecell{0.99} & \makecell{0.94} & \makecell{0.88} \\
\hline
\makecell{사다리} & \makecell{1.01} & \makecell{1.01} & \makecell{1.00} & \makecell{0.95} & \makecell{0.89} \\
\hline
\makecell{부정형} & \makecell{1.06} & \makecell{1.06} & \makecell{1.05} & \makecell{1.00} & \makecell{0.94} \\
\hline
\makecell{자루형} & \makecell{1.14} & \makecell{1.14} & \makecell{1.13} & \makecell{1.07} & \makecell{1.00} \\
\hline
\end{tabular}\par}

※ 장방형(가장형, 세장형) \par
{\small\centering
\begin{tabular}{|c|c|c|c|c|}
\hline
\makecell{공공용지} & \makecell{일반} & \makecell{도로} & \makecell{공원} & \makecell{운동장} \\
\hline
\makecell{일반} & \makecell{1.00} & \makecell{0.33} & \makecell{0.33} & \makecell{0.70} \\
\hline
\makecell{도로} & \makecell{3.03} & \makecell{1.00} & \makecell{1.00} & \makecell{2.12} \\
\hline
\makecell{공원} & \makecell{3.03} & \makecell{1.00} & \makecell{1.00} & \makecell{2.12} \\
\hline
\makecell{운동장} & \makecell{1.43} & \makecell{0.47} & \makecell{0.47} & \makecell{1.00} \\
\hline
\end{tabular}\par}

- H동(주거지역) \par
{\small\centering
\begin{tabular}{|c|c|c|c|c|}
\hline
\makecell{이용상황} & \makecell{주거용} & \makecell{상업업무} & \makecell{주상용} & \makecell{공업용} \\
\hline
\makecell{주거용} & \makecell{1.00} & \makecell{1.29} & \makecell{1.15} & \makecell{1.08} \\
\hline
\makecell{상업업무} & \makecell{0.78} & \makecell{1.00} & \makecell{0.89} & \makecell{0.84} \\
\hline
\makecell{주상용} & \makecell{0.87} & \makecell{1.12} & \makecell{1.00} & \makecell{0.94} \\
\hline
\makecell{공업용} & \makecell{0.93} & \makecell{1.19} & \makecell{1.06} & \makecell{1.00} \\
\hline
\end{tabular}\par}

{\small\centering
\begin{tabular}{|c|c|c|c|c|}
\hline
\makecell{도로접면} & \makecell{소로한면} & \makecell{소로각지} & \makecell{세로(가)} & \makecell{세각(가)} \\
\hline
\makecell{소로한면} & \makecell{1.00} & \makecell{1.04} & \makecell{0.94} & \makecell{0.99} \\
\hline
\makecell{소로각지} & \makecell{0.97} & \makecell{1.00} & \makecell{0.91} & \makecell{0.95} \\
\hline
\makecell{세로(가)} & \makecell{1.06} & \makecell{1.10} & \makecell{1.00} & \makecell{1.05} \\
\hline
\makecell{세각(가)} & \makecell{1.01} & \makecell{1.05} & \makecell{0.95} & \makecell{1.00} \\
\hline
\end{tabular}\par}

{\small\centering
\begin{tabular}{|c|c|c|c|c|c|}
\hline
\makecell{고저} & \makecell{저지} & \makecell{평지} & \makecell{완경사} & \makecell{급경사} & \makecell{고지} \\
\hline
\makecell{저지} & \makecell{1.00} & \makecell{1.10} & \makecell{0.99} & \makecell{0.90} & \makecell{0.89} \\
\hline
\makecell{평지} & \makecell{0.91} & \makecell{1.00} & \makecell{0.90} & \makecell{0.82} & \makecell{0.81} \\
\hline
\makecell{완경사} & \makecell{1.01} & \makecell{1.11} & \makecell{1.00} & \makecell{0.91} & \makecell{0.90} \\
\hline
\makecell{급경사} & \makecell{1.11} & \makecell{1.22} & \makecell{1.10} & \makecell{1.00} & \makecell{0.99} \\
\hline
\makecell{고지} & \makecell{1.12} & \makecell{1.24} & \makecell{1.11} & \makecell{1.01} & \makecell{1.00} \\
\hline
\end{tabular}\par}

{\small\centering
\begin{tabular}{|c|c|c|c|c|c|}
\hline
\makecell{형상(주거용)} & \makecell{정방형} & \makecell{장방형} & \makecell{사다리} & \makecell{부정형} & \makecell{자루형} \\
\hline
\makecell{정방형} & \makecell{1.00} & \makecell{1.01} & \makecell{0.96} & \makecell{0.93} & \makecell{0.87} \\
\hline
\makecell{장방형} & \makecell{0.99} & \makecell{1.00} & \makecell{0.95} & \makecell{0.92} & \makecell{0.86} \\
\hline
\makecell{사다리} & \makecell{1.04} & \makecell{1.05} & \makecell{1.00} & \makecell{0.97} & \makecell{0.91} \\
\hline
\makecell{부정형} & \makecell{1.08} & \makecell{1.09} & \makecell{1.03} & \makecell{1.00} & \makecell{0.94} \\
\hline
\makecell{자루형} & \makecell{1.15} & \makecell{1.16} & \makecell{1.10} & \makecell{1.07} & \makecell{1.00} \\
\hline
\end{tabular}\par}

※ 장방형(가장형, 세장형) \par
{\small\centering
\begin{tabular}{|c|c|c|c|c|}
\hline
\makecell{공공용지} & \makecell{일반} & \makecell{도로} & \makecell{공원} & \makecell{운동장} \\
\hline
\makecell{일반} & \makecell{1.00} & \makecell{0.33} & \makecell{0.33} & \makecell{0.70} \\
\hline
\makecell{도로} & \makecell{3.03} & \makecell{1.00} & \makecell{1.00} & \makecell{2.12} \\
\hline
\makecell{공원} & \makecell{3.03} & \makecell{1.00} & \makecell{1.00} & \makecell{2.12} \\
\hline
\makecell{운동장} & \makecell{1.43} & \makecell{0.47} & \makecell{0.47} & \makecell{1.00} \\
\hline
\end{tabular}\par}

< 자료 6 > 사례자료 \par
1. 평가사례 \par
{\small\centering\footnotesize\setlength{\tabcolsep}{3pt}
\begin{tabular}{|c|c|c|c|c|c|c|c|c|c|}
\hline
\makecell{기호} & \makecell{기준시점} & \makecell{소재지} & \makecell{면적\newline{}(m\textsuperscript{2})} & \makecell{지목} & \makecell{이용상황} & \makecell{용도\newline{}지역} & \makecell{도로\newline{}접면} & \makecell{형상\newline{}지세} & \makecell{토지단가\newline{}(원/m\textsuperscript{2})} \\
\hline
\multirow{2}{*}{\makecell{㉠}} & \makecell{2025-\newline{}12-23} & \makecell{M동\newline{}406-1} & \makecell{94.2} & \makecell{대} & \makecell{단독주택} & \makecell{2종일주} & \makecell{세로\newline{}(가)} & \makecell{세장형\newline{}평지} & \makecell{12,100,000} \\
\cline{2-10}
 & \multicolumn{9}{|c|}{\makecell{- 평가목적: 담보\newline{}- 기준가치: 시장가치}} \\
\hline
\multirow{2}{*}{\makecell{㉡}} & \makecell{2026-\newline{}01-07} & \makecell{M동\newline{}501-2} & \makecell{278.2} & \makecell{대} & \makecell{단독주택} & \makecell{2종일주} & \makecell{소로\newline{}각지} & \makecell{가장형\newline{}평지} & \makecell{13,300,000} \\
\cline{2-10}
 & \multicolumn{9}{|c|}{\makecell{- 평가목적: 담보\newline{}- 기준가치: 시장가치}} \\
\hline
\multirow{2}{*}{\makecell{㉢}} & \makecell{2025-\newline{}03-15} & \makecell{H동\newline{}447-1} & \makecell{141.2} & \makecell{대} & \makecell{단독주택} & \makecell{2종일주} & \makecell{세로\newline{}(가)} & \makecell{정방형\newline{}평지} & \makecell{13,310,000} \\
\cline{2-10}
 & \multicolumn{9}{|c|}{\makecell{- 평가목적: 소송(소유권이전 사건)\newline{}- ‘H동 일대 가로주택정비사업’구역 내 토지}} \\
\hline
\end{tabular}\par}

2. 거래사례 \par
{\small\centering\setlength{\tabcolsep}{3pt}\footnotesize\setlength{\tabcolsep}{3pt}
\begin{tabular}{|c|c|c|c|c|c|c|c|c|c|}
\hline
\makecell{기호} & \makecell{거래\newline{}시점} & \makecell{소재지} & \makecell{면적\newline{}(m\textsuperscript{2})} & \makecell{지목} & \makecell{이용상황} & \makecell{용도\newline{}지역} & \makecell{도로접면} & \makecell{형상\newline{}지세} & \makecell{토지단가\newline{}(원/m\textsuperscript{2})} \\
\hline
\multirow{2}{*}{\makecell{ⓐ}} & \makecell{2025-\newline{}11-02} & \makecell{M동\newline{}562-1} & \makecell{122.3} & \makecell{대} & \makecell{단독주택} & \makecell{2종\newline{}일주} & \makecell{세로\newline{}(가)} & \makecell{정방형\newline{}평지} & \makecell{12,100,000} \\
\cline{2-10}
 & \multicolumn{9}{|p{0.90\linewidth}|}{- 거래가액: 1,480,000,000원\newline{}- 경과년수, 현상 및 관리상태 등에 비추어 건물 가치는 없는 것으로 판단됨} \\
\hline
\multirow{2}{*}{\makecell{ⓑ}} & \makecell{2026-\newline{}03-23} & \makecell{M동\newline{}570-1} & \makecell{135.4} & \makecell{대} & \makecell{단독주택} & \makecell{2종\newline{}일주} & \makecell{세로\newline{}(가)} & \makecell{사다리\newline{}평지} & \makecell{11,300,000} \\
\cline{2-10}
 & \multicolumn{9}{|p{0.90\linewidth}|}{- 거래가액: 1,530,000,000원\newline{}- 경과년수, 현상 및 관리상태 등에 비추어 건물가치는 없는 것으로 판단됨\newline{}- ‘M동 일대 가로주택정비사업’구역 내 토지로서, 거래 신고는 되어 있으나\newline{}소유권이전등기는 완료되지 않은 상태임\newline{}- 거래가액이 인근지역 지가수준에 비하여 낮은 것으로 판단됨} \\
\hline
\end{tabular}\par}

※ 정비사업 조합설립인가 후에 소유권을 취득하는 경우, 원칙적으로 분양권이 주어지지 않는 등 정부의 규제조치로 서울특별시 내 분양권 거래는 현실적으로 이루어지지 않는 상태임 \par
< 자료 7 > ‘M동 일대 가로주택정비사업’ 분양신청 안내서 주요 내용 \par
1. 분양대상 대지 또는 건축물 내역 [단위: 세대, m\textsuperscript{2}] \par
{\small\centering
\begin{tabular}{|c|c|c|c|c|c|c|}
\hline
\multirow{2}{*}{\makecell{타입}} & \multirow{2}{*}{\makecell{세대수}} & \multicolumn{3}{|c|}{\makecell{공급면적}} & \multirow{2}{*}{\makecell{기타공용}} & \multirow{2}{*}{\makecell{계약면적}} \\
\cline{3-5}
 &  & \makecell{전유면적} & \makecell{공용면적} & \makecell{계} &  &  \\
\hline
\makecell{59} & \makecell{170} & \makecell{59.99} & \makecell{22.88} & \makecell{82.87} & \makecell{81.57} & \makecell{164.44} \\
\hline
\makecell{84} & \makecell{95} & \makecell{84.99} & \makecell{30.41} & \makecell{115.40} & \makecell{115.56} & \makecell{230.96} \\
\hline
\makecell{합계} & \makecell{265} & \makecell{18,272.35} & \makecell{6,778.55} & \makecell{25,050.9} & \makecell{24,845.10} & \makecell{49,896.00} \\
\hline
\end{tabular}\par}

2. 조합원 예상 분양가 \par
{\small\centering
\begin{tabular}{|c|c|c|c|}
\hline
\multicolumn{2}{|c|}{\makecell{단위세대}} & \multirow{2}{*}{\makecell{세대별 예상 분양가}} & \multirow{2}{*}{\makecell{비고}} \\
\cline{1-2}
\makecell{타입} & \makecell{공급면적(m\textsuperscript{2})} &  &  \\
\hline
\makecell{59} & \makecell{82.87} & \makecell{990,820,000원} & \makecell{-} \\
\hline
\makecell{84} & \makecell{115.40} & \makecell{1,379,800,000원} & \makecell{-} \\
\hline
\end{tabular}\par}

3. 비례율: 101.13\% \par
4. 본건 종전자산평가 내역 \par
{\small\centering
\begin{tabular}{|c|c|c|c|c|c|c|}
\hline
\multicolumn{2}{|c|}{\makecell{소재지}} & \makecell{토지면적(m\textsuperscript{2})} & \makecell{토지\newline{}평가액(원)} & \makecell{건물면적(m\textsuperscript{2})} & \makecell{건물\newline{}평가액(원)} & \makecell{합계금액(원)} \\
\hline
\makecell{지번} & \makecell{동호수} & \multirow{2}{*}{\makecell{108.0}} & \multirow{2}{*}{\makecell{431,244,000}} & \multirow{2}{*}{\makecell{-}} & \multirow{2}{*}{\makecell{-}} & \multirow{2}{*}{\makecell{431,244,000}} \\
\cline{1-2}
\makecell{567-12} & \makecell{-} &  &  &  &  &  \\
\hline
\end{tabular}\par}

※ ‘서울특별시 빈집 및 소규모주택 정비에 관한 조례’에 따르면 가로주택정비사업에서 종전토지의 총면적이 90m\textsuperscript{2} 이상인 경우 분양대상자가 됨 \par
5. 본건 분양주택의 예상 시장가치 \par
- 1) 예상 분양주택 타입 \par
본 가로주택정비사업의 예상 분양주택은 총 265세대로서, 이 중 59타입이 전체의 약 64\%를 차지하고 선호도가 높으며 본건 토지의 종전자산평가액이 상대적으로 낮은 점등을 고려하여 59타입을 예상 분양주택으로 선정함 \par
- 2) 59타입 분양주택 시장가치 추정 \par
- 본 사업구역 인근 유사한 단지 규모의 59타입 가격수준은 신축(입주)년도 등에 따라 14.5억원(약 10년 경과) ∼ 16.5억원(약 5년 경과)을 나타내고 있음 \par
- 최근 몇 년간 물가상승 등의 영향으로 건축비는 상승되어 왔으며, 서울특별시 전체적으로 아파트 임대가(전월세) 또한 지속적인 상승세를 나타내고 있음 \par
- 인근지역에 대한 시장조사 결과, 본건 정비사업의 예상 추가분담금 납부 만료일 기준으로 59타입 분양주택의 시장가치는 19억원(시장조사 범위값: 18억원 ∼ 20억원)으로 추정됨 \par
< 자료 8 > 시장금리 및 할인율 \par
1. 시장금리 추이 \par
[단위: 연, \%] \par
{\small\centering\footnotesize\setlength{\tabcolsep}{3pt}
\begin{tabular}{|c|c|c|c|c|c|c|c|c|}
\hline
\makecell{항목} & \makecell{2024\newline{}2Q} & \makecell{2024\newline{}3Q} & \makecell{2024\newline{}4Q} & \makecell{2025\newline{}1Q} & \makecell{2025\newline{}2Q} & \makecell{2025\newline{}3Q} & \makecell{2025\newline{}4Q} & \makecell{2026\newline{}1Q} \\
\hline
\makecell{국고채\newline{}(3년)} & \makecell{3.380} & \makecell{2.967} & \makecell{2.785} & \makecell{2.591} & \makecell{2.389} & \makecell{2.453} & \makecell{2.844} & \makecell{3.198} \\
\hline
\makecell{국고채\newline{}(5년)} & \makecell{3.419} & \makecell{2.998} & \makecell{2.853} & \makecell{2.698} & \makecell{2.519} & \makecell{2.608} & \makecell{3.029} & \makecell{3.458} \\
\hline
\makecell{회사채\newline{}(3년, AA-)} & \makecell{3.857} & \makecell{3.472} & \makecell{3.382} & \makecell{3.197} & \makecell{2.957} & \makecell{2.933} & \makecell{3.292} & \makecell{3.746} \\
\hline
\makecell{회사채\newline{}(3년, BBB-)} & \makecell{10.029} & \makecell{9.461} & \makecell{9.199} & \makecell{8.960} & \makecell{8.718} & \makecell{8.770} & \makecell{9.139} & \makecell{9.570} \\
\hline
\end{tabular}\par}

2. 할인율 등 \par
- 시장금리는 2025년 2분기 이후 전반적으로 상승세를 보이고 있으며, 물가상승이 지속될 경우 시장금리 또한 지속적으로 상승할 것으로 예상됨 \par
- 본건과 같은 주택정비사업은 다수인의 이해관계가 복잡하게 작용하여 사업기간이 연장될 수 있고, 이에 따라 추가분담금이 증가하는 등 미래에 대한 불확실성 등으로 투자 위험은 기존 아파트에 대한 투자에 비해 높은 것이 일반적임 \par
- 인근지역에 대한 시장조사 결과를 기준으로 본건 정비사업에 적용할 할인율은 6.5\%/연(범위값: 6.0\%/연 ∼ 7.0\%/연)로 결정함 \par
<자료 9> 본건 사업시행자 제시 정비사업 추진계획 \par
2026. 12.: 사업시행계획인가(관리처분계획 포함) \par
2027. 02.: 이주 개시 \par
2027. 12.: 착공 \par
2030. 12.: 준공 및 입주 \par
2031. 02.: 추가분담금 납부 기한 \par
<자료 10> 기타 사항 \par
1. 지역요인 및 개별요인 비교치는 소수점 셋째자리 이하 절사 \par
2. 그 밖의 요인 보정치는 표준지 기준 산정방식을 적용하고, 소수점 셋째자리 이하 절사 \par
3. 공시지가기준법에 의한 시산단가는 백원단위에서 절사하여 천원단위까지 표시함 \par
4. 개발법 적용 시 할인기간은 연간단위를 기준으로 적용하되, 10개월 이상인 경우 1년으로 간주함 \par
5. 각 감정평가방법에 의한 시산가액은 백원단위에서 절사하여 천원단위까지 표시함 \par
\end{longquestion}
